% GENERATED FILE. Edit canonical lessons, then run npm run generate:books.
% canonical-source-hash: fnv1a64:13d5639b9a122554
% canonical-lessons: JA-C282-sekai, JA-C282-sensou, JA-C282-heiwa, JA-C282-tatakau, JA-C282-seifu

\chapter{News from Abroad}
\label{ch:ja-news-from-abroad}
\hlchaptermodality{282}

\begin{chapteropening}
\textbf{By the end of this chapter:} \emph{I can say what is happening between countries.}

Complete the last lesson of chapter 282: I can say what is happening between countries.
\end{chapteropening}

\section[\texorpdfstring{\ja{せかい} (sekai)}{せかい (sekai)}]{\ja{せかい} (sekai) — the world}

\label{lesson:JA-C282-sekai}



\begin{quote}
Before the new one: say the Japanese for a topic of conversation, then the Japanese for an uproar.
\end{quote}

\subsection*{You'll want to know: \ja{せかい}}
\textbf{\ja{せかい}} — \emph{sekai} — \textquotedblleft{}the world\textquotedblright{}.

\begin{quote}
\textbf{\ja{せかいで} \ja{いちばん} \ja{たかい} \ja{やまは} \ja{どこですか。}} — \emph{sekai de ichiban takai yama wa doko desu ka} — \textquotedblleft{}Which is the highest mountain in the world?\textquotedblright{}
\end{quote}

\begin{cousinweb}
\textbf{\ja{せかい}} (\emph{sekai}) is Sino-Japanese (kango): two Chinese characters read with their on'yomi, readings borrowed from Chinese. They are \emph{se}, \textquotedblleft{}world\textquotedblright{} (Middle Chinese \emph{syejH}) and \emph{kai}, \textquotedblleft{}boundary\textquotedblright{} (\emph{keajH}) (KANJIDIC2; Middle Chinese in Baxter's transcription).
\end{cousinweb}

\begin{grammarlens}[title={What you've built}]
One more word for telling what happened.
\end{grammarlens}

\subsection*{Your turn}
\begin{itemize}
\raggedright
  \item \emph{Say it:} \emph{sekai}
  \item \emph{Say it:} \emph{sekai}, once more
  \item \emph{Say it:} \emph{sekai de ichiban takai yama wa doko desu ka} — Which is the highest mountain in the world?
  \item \emph{Recall:} say the Japanese for a topic of conversation, then the Japanese for an uproar, then say \emph{sekai} again
\end{itemize}

\begin{tcolorbox}[breakable,colback=teal!4,colframe=teal!35!black,title={Before you move on}]
What does \ja{せかい} mean? (\textquotedblleft{}The world\textquotedblright{}.) Say it once more.
\end{tcolorbox}
\section[\texorpdfstring{\ja{せんそう} (sensō)}{せんそう (sensō)}]{\ja{せんそう} (sensō) — a war}

\label{lesson:JA-C282-sensou}



\begin{quote}
Before the new one: say the Japanese for an uproar, then the Japanese for the world.
\end{quote}

\subsection*{You'll want to know: \ja{せんそう}}
\textbf{\ja{せんそう}} — \emph{sensō} — \textquotedblleft{}a war\textquotedblright{}.

\begin{quote}
\textbf{\ja{その} \ja{くには} \ja{せんそうを} \ja{しています。}} — \emph{sono kuni wa sensō o shite imasu} — \textquotedblleft{}That country is at war.\textquotedblright{}
\end{quote}

\begin{cousinweb}
\textbf{\ja{せんそう}} (\emph{sensō}) is Sino-Japanese (kango): two Chinese characters read with their on'yomi, readings borrowed from Chinese. They are \emph{sen}, \textquotedblleft{}war\textquotedblright{} (Middle Chinese \emph{tsyenH}) and \emph{sō}, \textquotedblleft{}contend\textquotedblright{} (\emph{tsreang}) (KANJIDIC2; Middle Chinese in Baxter's transcription).
\end{cousinweb}

\begin{grammarlens}[title={What you've built}]
One more word for telling what happened.
\end{grammarlens}

\subsection*{Your turn}
\begin{itemize}
\raggedright
  \item \emph{Say it:} \emph{sensō}
  \item \emph{Say it:} \emph{sensō}, once more
  \item \emph{Say it:} \emph{sono kuni wa sensō o shite imasu} — That country is at war.
  \item \emph{Recall:} say the Japanese for an uproar, then the Japanese for the world, then say \emph{sensō} again
\end{itemize}

\begin{tcolorbox}[breakable,colback=teal!4,colframe=teal!35!black,title={Before you move on}]
What does \ja{せんそう} mean? (\textquotedblleft{}A war\textquotedblright{}.) Say it once more.
\end{tcolorbox}
\section[\texorpdfstring{\ja{へいわ} (heiwa)}{へいわ (heiwa)}]{\ja{へいわ} (heiwa) — peace}

\label{lesson:JA-C282-heiwa}



\begin{quote}
Before the new one: say the Japanese for the world, then the Japanese for a war.
\end{quote}

\subsection*{You'll want to know: \ja{へいわ}}
\textbf{\ja{へいわ}} — \emph{heiwa} — \textquotedblleft{}peace\textquotedblright{}.

\begin{quote}
\textbf{\ja{みんな} \ja{へいわを} \ja{ねがっています。}} — \emph{minna heiwa o negatte imasu} — \textquotedblleft{}Everyone wishes for peace.\textquotedblright{}
\end{quote}

A na-adjective: \ja{な} before a noun.

\begin{cousinweb}
\textbf{\ja{へいわ}} (\emph{heiwa}) is Sino-Japanese (kango): two Chinese characters read with their on'yomi, readings borrowed from Chinese. They are \emph{hei}, \textquotedblleft{}peace\textquotedblright{} (Middle Chinese \emph{bjaeng}) and \emph{wa}, \textquotedblleft{}harmony\textquotedblright{} (\emph{hwa}) (KANJIDIC2; Middle Chinese in Baxter's transcription).
\end{cousinweb}

\begin{grammarlens}[title={What you've built}]
One more word for telling what happened.
\end{grammarlens}

\subsection*{Your turn}
\begin{itemize}
\raggedright
  \item \emph{Say it:} \emph{heiwa}
  \item \emph{Say it:} \emph{heiwa}, once more
  \item \emph{Say it:} \emph{minna heiwa o negatte imasu} — Everyone wishes for peace.
  \item \emph{Recall:} say the Japanese for the world, then the Japanese for a war, then say \emph{heiwa} again
\end{itemize}

\begin{tcolorbox}[breakable,colback=teal!4,colframe=teal!35!black,title={Before you move on}]
What does \ja{へいわ} mean? (\textquotedblleft{}Peace\textquotedblright{}.) Say it once more.
\end{tcolorbox}
\section[\texorpdfstring{\ja{たたかう} (tatakau)}{たたかう (tatakau)}]{\ja{たたかう} (tatakau) — to fight; to battle}

\label{lesson:JA-C282-tatakau}



\begin{quote}
Before the new one: say the Japanese for a war, then the Japanese for peace.
\end{quote}

\subsection*{You'll want to know: \ja{たたかう}}
\textbf{\ja{たたかう}} — \emph{tatakau} — \textquotedblleft{}to fight; to battle\textquotedblright{}.

\begin{quote}
\textbf{\ja{ふたつの} \ja{くにが} \ja{たたかいました。}} — \emph{futatsu no kuni ga tatakaimashita} — \textquotedblleft{}The two countries fought each other.\textquotedblright{}
\end{quote}

Polite \textbf{\ja{たたかいます}}, \emph{tatakaimasu}; a godan verb. Past (polite): \textbf{\ja{その} \ja{くにが} \ja{たたかいました。}}, \emph{sono kuni ga tatakaimashita}.

\begin{cousinweb}
\textbf{\ja{たたかう}} (\emph{tatakau}) is a native Japanese word (wago), inherited rather than borrowed from Chinese; its kanji was attached later as a kun'yomi (JMdict).
\end{cousinweb}

\begin{grammarlens}[title={What you've built}]
One more word for telling what happened.
\end{grammarlens}

\subsection*{Your turn}
\begin{itemize}
\raggedright
  \item \emph{Say it:} \emph{tatakau}
  \item \emph{Say it:} \emph{tatakau}, once more
  \item \emph{Say it:} \emph{futatsu no kuni ga tatakaimashita} — The two countries fought each other.
  \item \emph{Recall:} say the Japanese for a war, then the Japanese for peace, then say \emph{tatakau} again
\end{itemize}

\begin{tcolorbox}[breakable,colback=teal!4,colframe=teal!35!black,title={Before you move on}]
What does \ja{たたかう} mean? (\textquotedblleft{}To fight; to battle\textquotedblright{}.) Say it once more.
\end{tcolorbox}
\section[\texorpdfstring{\ja{せいふ} (seifu)}{せいふ (seifu)}]{\ja{せいふ} (seifu) — the government}

\label{lesson:JA-C282-seifu}



\begin{quote}
Before the new one: say the Japanese for peace, then the Japanese for to fight.
\end{quote}

\subsection*{You'll want to know: \ja{せいふ}}
\textbf{\ja{せいふ}} — \emph{seifu} — \textquotedblleft{}the government\textquotedblright{}.

\begin{quote}
\textbf{\ja{せいふは} \ja{なにも} \ja{いいませんでした。}} — \emph{seifu wa nanimo iimasen deshita} — \textquotedblleft{}The government said nothing.\textquotedblright{}
\end{quote}

\begin{cousinweb}
\textbf{\ja{せいふ}} (\emph{seifu}) is Sino-Japanese (kango): two Chinese characters read with their on'yomi, readings borrowed from Chinese. They are \emph{sei}, \textquotedblleft{}government\textquotedblright{} (Middle Chinese \emph{tsyengH}) and \emph{fu}, \textquotedblleft{}govt office\textquotedblright{} (\emph{pjuX}) (KANJIDIC2; Middle Chinese in Baxter's transcription).
\end{cousinweb}

\begin{grammarlens}[title={What you've built}]
One more word for telling what happened.
\end{grammarlens}

\subsection*{Your turn}
\begin{itemize}
\raggedright
  \item \emph{Say it:} \emph{seifu}
  \item \emph{Say it:} \emph{seifu}, once more
  \item \emph{Say it:} \emph{seifu wa nanimo iimasen deshita} — The government said nothing.
  \item \emph{Recall:} say the Japanese for peace, then the Japanese for to fight, then say \emph{seifu} again
\end{itemize}

\begin{tcolorbox}[breakable,colback=teal!4,colframe=teal!35!black,title={Before you move on}]
What does \ja{せいふ} mean? (\textquotedblleft{}The government\textquotedblright{}.) Say it once more.
\end{tcolorbox}
